\chapter{LITERATURE REVIEW}
\section{Rework as an Evidence-Dependent Construction Event}
\textcite{forcada2017} examine the cost of distinguishable rework incidents; \textcite{liu2020} present building-project case insights; and \textcite{love2022} review how errors, organizational conditions and interventions can generate rework and its consequences. These studies justify a case-based approach rather than an assumption that all observed defects involved paid corrective work. \textcite{asadi2023} provide a classification perspective for rework causes and potentially liable participants. The present study deliberately separates the description of a construction condition, a corrective mechanism, supported causation and documentary evidence of resource impacts.

Construction-management literature also points to project control, coordination and changes. \textcite{love2019} examine rework and changes in relation to contractor margins, while \textcite{safapour2019} analyse manageable causes and prevention. \textcite{hwang2019} consider information-management practices, and \textcite{garg2021} investigate cause relationships in a developing-country setting. \textcite{asadi2021} review contracting-strategy contexts; \textcite{mostofi2022} address high-cost rework prediction. These ten papers constitute the previously established scholarly basis of the manuscript. They are not substitutes for records from an actual site in Pokhara.

\section{Waterproofing as a Construction System}
A tiled bathroom is an assembly of floor and wall substrates, drainage falls, membrane or coating, corners and transitions, penetrations, adhesives and finishes. Tile and grout should not be assumed to constitute a complete waterproof barrier; a dedicated system is necessary where specified for sustained wet service \parencite{schluter2026,tcna2023}. In a wet-area assembly, the membrane's continuity at corners, walls, door thresholds and drain interfaces is at least as significant as field-area application. Manufacturer-specific products and substrate compatibility govern details such as layer thickness, sealing accessories and the required height of vertical protection \parencite{laticrete2026}.

Omitted membranes can permit migration to adjoining finishes when water is present and a flow path exists; however, the visible location of staining does not independently reveal the source of the water. Leakage diagnosis requires inspection of the plumbing, drain system, wet-area assembly and exterior exposure before the probable route is established. At roof terraces, deficient drainage falls may increase ponding and exposure, while failed laps, upstands, terminations, penetrations or protective layers may also create an entry path. Membrane compatibility with cement-based plaster or tile adhesive varies: a smooth surface is not proof that every sheet membrane lacks adhesion, and a manufacturer's approved bonded system must be consulted.

\section{Masonry Interfaces, Mortar and Plaster}
Cracking near intersections of walls can be associated with differential movement, discontinuity of the masonry bond, inadequate connection details or incompatible movement of adjoining elements. A vertical butt joint is not equivalent to interlocking bonded masonry; engineered ties or other approved connection details can be appropriate alternatives when correctly designed and installed. A visual line in plaster does not establish an inadequately bonded wall without exposing or documenting the substrate.

The initial rate of water absorption of clay bricks influences their withdrawal of moisture from fresh mortar, making brick condition and mortar water retention relevant to bond development. \textcite{bailey1990} discuss the significance of brick suction; technical recommendations for fired-clay units depend on material properties and construction conditions \parencite{bia2026}. Accordingly, universal soaking instructions are inappropriate. Inadequate moisture control during subsequent mortar curing is a separate process concern. Neither mortar strength nor the cause of a particular crack can be established without proper records and, where needed, testing.

Nepal's Department of Urban Development and Building Construction (DUDBC) publishes NBC 205:2024, a ready-to-use guideline expressly framed for low-rise reinforced concrete buildings \emph{without masonry infill}, and also lists NBC 201 for reinforced concrete buildings with masonry infill \parencite{dudbc2024}. Details on bands and non-load-bearing wall reinforcement must be applied according to the building's actual structural system, relevant provisions and approved drawings. The research therefore treats reports of omitted lintel and sill bands as matters for drawing and code compliance checks, not automatic proof of a breach in every building or of the physical cause of each crack.

\section{Integrated Analytical Framework}
The same evidence chain is used across the two building-work subjects: (1) field concern, (2) physical condition and location, (3) plausible but testable mechanism, (4) requirement from the approved design or applicable technical reference, (5) evidence of original executed work, (6) evidence of corrective action, and (7) supported resource or schedule consequence. Causes and consequences are not imported automatically from international studies. The framework allows a practitioner observation to form the basis for discussion without converting it into a verified site-case frequency.

\begin{table}[htbp]
\caption{Distinguishing the evidence levels used in this thesis}
\label{tab:evidencelevels}
\centering
\begin{tabularx}{\textwidth}{@{}L{3.6cm}X@{}}
\toprule Evidence level & What the level establishes \\ \midrule
Field observation & A visible or witnessed construction condition is described; available records may not establish its cause or subsequent correction.\\
Technical interpretation & A plausible mechanism explained by engineering principles or a published source; it is not a field-proven cause.\\
Verified corrective work & Records identify prior work and an actual subsequent intervention at a traceable place and time.\\
Attributed impact & Supporting quantities, payroll, material records or programme evidence permit consequences to be assessed without double counting.\\ \bottomrule
\end{tabularx}
\end{table}

\section{Research Gap}
The established ten-paper rework literature provides general causal and management perspectives, while the technical references explain specific workmanship and materials mechanisms. Neither body of literature determines the extent, frequency or cost of the two building-work concerns examined in Pokhara. A valuable local contribution at the present evidence level is a technically disciplined synthesis of reported concerns and a practical method to test and document them. A quantified, site-level rework comparison remains a future evidentiary task.
